% Research candidate only; no main-manuscript edit.
% The allowed cycles are the existing D_E,D_M,D_A and all their images.
% This is not a claim about arbitrary Hamilton cycles of the fibres.
\paragraph{A cycle-deletion bound.}
Let $C,D$ be cycles, let $U=V(C)\cap V(D)\ne\varnothing$, and let
$H$ contain all their edges with both endpoints in $U$. Let $F\subseteq
E(H)$ be forbidden, and let $R\subseteq E(H)\setminus F$ be used by
neither cycle. If the vertices outside $U$ form at most $r_C,r_D$
paths in the respective cycles, counting isolated vertices as paths,
then
\[
 |(E(C)\cap E(D))\setminus F|
 \ge 2|U|-r_C-r_D-|E(H)|-|F|+|R|.                 \tag{DP}
\]
Indeed, if the outside vertices form $q$ paths, deleting them leaves
exactly $|U|-q$ cycle edges; use $q=0$ for an empty outside set.
Removing $F$ loses at most $|F|$ edges from either cycle. Both remaining
edge sets lie in $E(H)\setminus(F\cup R)$, so inclusion--exclusion gives
the bound.

\paragraph{Deletion descriptions of the two unambiguous families.}
Let $a,b,c,d$ be distinct. For the end-change cycle at $aab$, deleting
the assignments with $w=c$ or $v=c$ removes one $P_4$; deleting those
with $u=c$ removes $P_3\cup P_2$. For the all-distinct cycle at $abc$,
the classes $v=a$, $u=b$, and $w=d$ respectively induce $P_6$,
$P_4\cup P_3$, and $P_3\cup P_2$ in the cycle. These are path
identities, not claims about induced paths in the whole fibre.
For a direct verification from \eqref{eq:path232-cycles}, their paths are
\[
\begin{array}{c|l}
E:\ w=c &[dbc,cbc,cdc,bdc]\\
E:\ v=c &[dca,bca,bcd,dcd]\\
E:\ u=c &[cbd,cba,cda],\ [cbc,cdc]\\[2pt]
A:\ v=a &[bab,bad,cad,dad,dab,cab]\\
A:\ u=b &[bad,bab,bdb,bcb],\ [bcd,bca,bda]\\
A:\ w=d &[bad,cad,dad],\ [dcd,bcd].
\end{array}                                                     \tag{DC}
\]
Here $[x_0,\ldots,x_k]$ denotes a path with these consecutive vertices.
The end-change special edges are $bcd$--$dcd$ and $bdc$--$cdc$.
Unused-colour interchange fixes the $E$ edge set, and reversal followed
by interchange of the endpoint boundary colours fixes the $A$ edge set.
Thus each fixed boundary has only one allowed $E$ or $A$ edge set,
and (DC) holds for every allowed cycle.

\paragraph{Three interfaces from the same count.}
Apply (DP) to
\[
 aab\to aac,\qquad aab\to acb,\qquad aab\to cab,
\]
with $H$ the common recolouring graph and
$F=(P_t\cup P_s)\cap E(H)$. Their fixed-middle layers are respectively
$(C_4,C_4,K_2)$, $(K_2\square K_3,C_4)$, and $(K_3,C_4,K_2)$.
There are respectively four, two, and three edges changing the middle
coordinate. Hence $(|U|,|E(H)|,|F|)$ is respectively
$(10,13,2)$, $(10,15,1)$, and $(9,11,1)$.
The special-edge counts follow from the two edges displayed above:
in the first interface one edge of each special set survives, and in
either other interface only one source special edge survives.
The deletion descriptions give $(r_C,r_D)=(1,1),(1,1),(2,2)$.
With $R=\varnothing$, (DP) gives respectively
\[
20-2-13-2=3,\qquad 20-2-15-1=2,\qquad 18-4-11-1=2
\]
safe common edges. Colour permutations and reversal give every labelled
interface of these three families.

\paragraph{Optional fourth family.}
For $abc\to abd$, the common graph has layers $(K_3,C_4,C_4)$ and
six middle-change edges, so $|U|=11$ and $|E(H)|=17$; $F$ is empty.
Deleting the respective unused endpoint colour gives two paths in
either cycle by (DC). The middle-$a$ layer is the triangle on
$bab,cab,dab$. Both cycles omit its edges $bab$--$cab$ and
$bab$--$dab$: at $bab$ the middle-$a$ path starts by changing the
last coordinate to the respective unused colour. Take these two edges
as $R$ in (DP). The lower bound becomes $22-4-17+2=3$.

% Scope: DP is a general cycle lemma. DC and the common-grid descriptions
% use the existing selected E/A cycles. No selected cycle is changed.
% Editorial comparison: the three primary families remove three entries
% of the present ten-entry table, and the optional fourth removes one
% more. Combining them with the separately frozen EE-transfer and
% A-middle support arguments leaves four entries in three rows: two
% EM orientations and the two MA source-normalization orbits. This note
% is longer than the current table and is not recommended for a pure
% page-count saving. EM and MA are not implied by the coarse bound.
